\documentclass{abntexto-ufc}

\ufcsetup{
  type = research-project,
  print-mode = single-sided,
  cover = true,
  institution = {Instituição de Origem Teste},
  submission-entity = {Universidade Federal do Ceará},
  project-program = {Programa de Pós-Graduação de Teste},
  project-type = {Projeto de pesquisa},
  author = {Autor Projeto Teste},
  title = {Projeto de Pesquisa Normativo},
  subtitle = {Validação da NBR 15287},
  title-variant = {Project Research Normative Test},
  volume = {2},
  location = {Fortaleza},
  year = {2026},
  advisor = {Prof. Orientador Projeto Teste}
}

\ufcAddBibliographyResource{tests/fixtures/references.bib}

\begin{document}
\pretextual

\ufcPrintCover
\ufcPrintTitlePage
\ufcPrintTableOfContents

\textual
\section{Introdução}
A introdução apresenta o tema, o problema, os objetivos e a justificativa do projeto.

\subsection{Problema de pesquisa}
Problema de pesquisa da fixture normativa.

\subsection{Objetivos}
Objetivos da fixture normativa.

\subsection{Justificativa}
Justificativa da fixture normativa.

\section{Referencial teórico}
Referencial teórico utilizado para embasar a proposta.

\section{Metodologia}
Métodos, procedimentos de coleta e estratégia de análise.

\section{Recursos}
Recursos necessários à execução da pesquisa.

\section{Cronograma}
Cronograma das atividades previstas para a pesquisa.

\nocite{silva2020}
\ufcPrintReferences

\end{document}
